\documentclass{article}

\usepackage{lipsum}
\usepackage{xcolor}
\usepackage[thmset=rich,
			sepcounters=true,
			proofenv=true,
			qedsymbolblackaquare=true,
			prooftitleitbf=true]{phfthm}
\usepackage{mdframed}

\phfMakeTheorem[defstar=true,
				thmstyle=plain]{answer}{Answer}

\newmdenv[subtitlebelowline=true,
		  subtitleaboveline=true,
		  subtitlebackgroundcolor=green,
		  backgroundcolor=blue,
		  frametitle={Custom mdframed Environment},
		  frametitlerule=true,
		  frametitlebackgroundcolor=yellow]
		 {mdfCustomed}

% New mdframed Theorem Environment
\mdfdefinestyle{mdfThmStyle}{linecolor=blue,
							 linewidth=2pt,
							 frametitlerule=true,
				    frametitlebackgroundcolor=gray!15}
\mdtheorem[style=mdfThmStyle]{mdfthm}{mdfThm} 

\begin{document}
	{ \centering
	  \Huge
	  \textbf{Theorems and Proofs Part II}
	}
	
	\vspace{2cm}
	
	\section{The phfthm Package}
	\begin{theorem}
		The Content to be Put here !!
	\end{theorem}
	\begin{question}
		The Content to be Put here !!
	\end{question}
	\begin{corollary*}
		The Content to be Put here !!
	\end{corollary*} 
	
	% A New User-defined Theorem-Like Environment
	\begin{thmheading}{Theorem-Like}
		\begin{equation}
			F(x) = 2x + y
		\end{equation}
	\end{thmheading}
	
	% A Typical Functioning User-defined Environment.
	\begin{answer}
		This is an Example of a Typical User-defined 			Theorem-Like Environment using phfthm Package.
	\end{answer}
	
	\begin{proof}
		This is a Proof Through the phfthm Package !!
	\end{proof}
	
	\pagebreak
	
	\section{The mdframed Package}
	\mdfdefinestyle{mymdfStyle}{linecolor=blue,
								linewidth=2pt}
								
	\begin{mdframed}[style=mymdfStyle]
		\lipsum[1]
	\end{mdframed}
	
	{
	\mdfapptodefinestyle{mymdfStyle}											{
						innertopmargin=1em,
						innerbottommargin=1em,
						innerleftmargin=1em,
						innerrightmargin=1em,
						fontcolor=cyan,
						backgroundcolor=gray!25,
						topline=false,
						bottomline=false,
						leftline=false}
						
	\begin{mdframed}[style=mymdfStyle,
					 frametitle={Lorem Ipsum Text:},
					 frametitlerule=true,
					 frametitlerulewidth=2pt,
					 frametitleaboveskip=0.4cm,
					 frametitlebelowskip=0.4cm]
		\lipsum[1]
	\end{mdframed}	
	}
	
	\begin{mdframed}[style=mymdfStyle]
		\lipsum[1-2]
	\end{mdframed}
	
	\begin{mdfCustomed}
		The Content \ldots
		\mdfsubtitle{The Subtitle}
		The Content \ldots
	\end{mdfCustomed}
	
	% A New User-defined Theorem Env. through mdframed.
	\begin{mdfthm*}[New Definition]
		Let $H$ be a subgroup of a group~$G$. A\emph{left coset} of $H$ in $G$ is a subset of $G$ that is of the form $xH$, where $x \in G$ and $xH = \{ xh : h \in H \}$. Similarly a \emph{right coset} of $H$ in $G$ is a subset of $G$ that is of the form $Hx$, where $Hx = \{ hx : h \in H \}$.	
	\end{mdfthm*}
	
\end{document}